\documentclass{article}
\usepackage{hyperref}
\title{ARIAC quality checks need a fresh observation after recovery}
\date{}
\begin{document}
\maketitle
This is a source-level calibration for I1, not an integrated factory
trial or a comparison with HeRo. It uses
\href{https://github.com/usnistgov/ARIAC/commit/ffc825757dc05cc6e4dfbd5c2483e993aab35652}{NIST ARIAC 2024.3.0}.
The pinned version is an independent test environment candidate;
it is not identified as the one used in the HeRo study.

The simulator's quality service reads a cached tray camera message
and scores the requested kit. In a fault-challenge quadrant, the
first check can also rename the actual simulated part to mark it
faulty and record that the quadrant was checked. A subsequent
check consults the part name retained in its chosen camera image.
The first call can therefore change the world without updating the
image used by the second. The response carries no camera sequence
or capture time. The camera's Gazebo transport callback and the
ROS service are separate in the source; a new request does not
itself require a new camera callback.

\paragraph{Executed check.}
The test compiles the original three functions
\texttt{PerformQualityCheck}, \texttt{HandleFaultyParts} and
\texttt{CheckQuadrantQuality} against a narrow fixture for the
world model, typed kit, cached image and message classes.
It uses a correct-type, correct-color part placed in the requested
quadrant for every branch. At the first check the live model
gets marked and the response indicates failure. Repeating the
check with the identical cached image gives a pass even though
the live model is still marked and has not been replaced. With
a new capture of the unchanged part the check fails again.
A capture after inserting a replacement passes. As a separate
control, an otherwise correct kit on the wrong tray has both
\texttt{all\_passed} and \texttt{incorrect\_tray} set. The client
must consume both. The last case changes only the tray identifier.

This is not a Gazebo reproduction. Camera interleaving and
world-image correspondence are inputs to the harness, not measured
delays or execution traces from an installed simulator. The source
bindings use the simulator's parsing convention for the part name;
physics, ROS queues, locking, real contacts and camera rate are
outside the harness. The replay never enters a missing-required-part
branch or models unsupported responses. A physical sensor blackout
need not stop this particular internal tray camera; that relationship
has not been tested and must not be inferred from generic sensor
health controls.

\paragraph{Consequences for a recovery boundary.}
A new service response is not proof of a new image or removal of
the old effect. An unchanged response can come from the stale image
or a real replacement. A useful readiness supplier would tie the
specific admission to a capture after the pertinent physical change
or exposure, verify the tray and required contents, and guard
against intervening changes before the operation consumes that
fact. In particular, a sample after the request is not enough if it
was taken before the removal it purports to verify. That conditional
design is an interface requirement, not an implementation using
the original response alone. Controller-visible authority, camera
provenance and actual completion must be supplied explicitly rather
than reading the simulator's ground truth.

Robot health names a separate condition. In the same ARIAC release
suction-control success changes the enable state, while actual joint
detachment is performed by the gripper update callback. The gripper
status contains an attachment bit and gripper type, but does not
name the dropped object or where it ended up. These statements come
from static source inspection only. Neither a resumed robot nor a
detached gripper certifies successful placement.

\paragraph{Published comparison scope.}
\href{https://arxiv.org/html/2103.11470\#S11}{NeBula} documents
separate mission watchdog, scheduled re-planning and heartbeat-regain
behaviors. The heartbeat behavior monitors renewed receipt; it does
not by itself supply the kit-effect history required here. The
\href{https://is.cas.cn/ztzl/rjsxsnh/2025xsnh/205hbzs/202607/P020260715632339212745.pdf}{authors' HeRo poster}
indicates pre-, during- and post-task monitors, synchronization and a
PDDL-fed recovery state machine. It does not license conclusions
about how their actual controller checks freshness, physical handoff
or concurrent custody. No whole-system test or superiority claim
is made. A matched direct controller with equivalent access to a
verified image can consume the same check.
\end{document}
